% !TeX root = ../教材题库_全部条目_转写工作稿.tex
% Source: XJ-B1-C04-EX-P095-Y0782
% 第4章 / 4.1 实数指数幂和幂函数 / 教材 p.95、96 / PDF 第102、103页
% OCR未校对：公式与文字均可能有误；未完成数学TeX转写；禁止直接发给学生
% 原题及相邻原解截图（核对用）：../题目原图/XJ-B1-C04-EX-P095-Y0782-01.webp, ../题目原图/XJ-B1-C04-EX-P095-Y0782-02.webp
% 本文件只在转写完成后作为题干使用；OCR 条目尚不是可靠题干。
\begin{bankitem}{4.1 例题 例题1}{第4章 / 4.1 实数指数幂和幂函数 / 教材 p.95、96 / PDF 第102、103页}{XJ-B1-C04-EX-P095-Y0782}{OCR未校对：公式与文字均可能有误；未完成数学TeX转写；禁止直接发给学生}
\begin{ocrdraft}
[教材p.95 / PDF第102页]
\par
00化简下列各式：
\par
（1）
\par
（2）
\par
（3）
\par
（3一a)
\par
（4）
\par
（5）
\par
（3一a)
\par
解
\par
（D（一2）3
\par
（2）
\par
（一2）4
\par
（3）
\par
（3一a）3
\par
（4）
\par
因为<，所以
\par
伍一仞2=
\par
磊；
\par
3一a，当磊<3时，
\par
（3一a）2
\par
a一3，当磊>3时．
\par
分数指数幂
\par
根式运算是一件比较复杂的事，例如，常常要先把根式化为同次根式再按运算
\par
法则进行运算，而引人分数指数幂的概念就可以大大简化根式运算．
\par
当a>0，m，7N且刀冫2时，规定
\par
这样就有
\par
2=2一
\par
，方便多了．
\par
[教材p.96 / PDF第103页]
\par
如果再规定0的正分数指数幂为0，0没有负分数指数幂，那么，对于任意有理
\par
数。、仍有下列运算法则：
\par
（）、=a气a>0，0s\ensuremath{\text{[OCR:EUR]}}Q)，
\par
Q=方．（a>0，b>0，reQ).
\par
这就把整数指数幂推广为有理数指数幂了．
\end{ocrdraft}
\end{bankitem}
